\SourceLeaf{404}%
\par
If we expand these two cosines, divide by \(\cos\tfrac12\varphi\cos\tfrac12\varphi'\), and denote \(\operatorname{tang}\tfrac12\varphi\) .by \(t\) and \(\operatorname{tang}\tfrac12\varphi'\) by \(t'\)
\[
\text{and}\quad\frac{a\sin A}{r-a\cos A}=\alpha,\quad\frac{r+a\cos A}{r-a\cos A}=\beta,\quad\text{we obtain}
\]
\[
\text{I.}\qquad t=\tfrac{1-\alpha t'}{\alpha-\beta t'}
\]
\par
In much the same way, if we set
\[
\operatorname{tang}\tfrac12\varphi''=t''\quad\frac{a'\sin A'}{r'-a'\cos A'}=\alpha',\quad\frac{r+a'\cos A'}{r-a'\cos A'}=\beta'
\]
we obtain
\[
\text{II.}\qquad t'=\tfrac{1-\alpha't''}{\alpha'-\beta't''}\ \text{etc.}
\]
\par
We see from this that we get as many equations as the polygon has sides, and that combining them ultimately yields a quadratic equation for \(t\).
\par\medskip\Needspace{5\baselineskip}
\pdfbookmark[2]{VI.}{carnot-zusatz-vi}
\centerline{\textbf{VI.}}\nobreak\smallskip
\textsc{Problem}. Given three circles in position and size, draw a fourth circle tangent to all three.

\textit{Solution}. Draw two mutually perpendicular axes through the centre of one circle and denote the distances from these lines by
\[
\begin{array}{rl}
\text{of the centre of the second circle}\quad .\ .\ .\ .\ .&a,\quad b\\
\text{of the third circle}\quad .\ .\ .\ .\ .&a',\quad b'\\
\text{of the required circle}\quad .\ .\ .\ .\ .&x,\quad y
\end{array}
\]
\par
Let the distance from the centre of the first circle to the centre of the required circle be \(=z\). Then
\begin{list}{}{\setlength{\leftmargin}{3em}\setlength{\labelwidth}{2.5em}\setlength{\labelsep}{.5em}}
\item[\(z-c\)] is the distance from the centre of the second circle to the centre of the required circle,
\item[\(z-c'\)] is the distance from the centre of the third circle to the centre of the required circle,
\end{list}
